\vspace{2.5pt plus 3pt minus 0.5pt}\noindent\hspace*{1.2em}\kbvnum{{[}2{]}}\hspace{0.55em}ಕರ್ತರಿ ಪ್ರಯೋಗದಲ್ಲಿ ಕ್ರಿಯಾಪದಗಳಿಗೆ ಕರ್ತೃವಿನ ಲಿಂಗವಚನಗಳೂ, ಕರ್ಮಣಿ ಪ್ರಯೋಗದಲ್ಲಿ ಕರ್ಮಪದದ ಲಿಂಗವಚನಗಳೂ ಬರುತ್ತವೆ.\par

\noindent\hspace*{1.2em}\kbvlabel{ಉದಾ}\hspace{0.6em}ಸುಜಾತ ಹಣ್ಣನ್ನು ತಿಂದಳು.\par

\noindent\hspace*{2.2em}ಚಿತ್ರವು ರವಿಯಿಂದ ಬರೆಯಲ್ಪಟ್ಟಿತು.\par

\vspace{2.5pt plus 3pt minus 0.5pt}\noindent\hspace*{1.2em}\kbvnum{{[}3{]}}\hspace{0.55em}ಒಂದೇ ಕ್ರಿಯೆಗೆ ವಿಭಿನ್ನ ಲಿಂಗದ ಅನೇಕ ಕರ್ತೃಪದಗಳಿರುವಾಗ ಕಟ್ಟಕಡೆಯ ಕರ್ತೃಪದದ ಲಿಂಗವೇ ಕ್ರಿಯಾಪದಕ್ಕೆ ಬರುವುದು.\par

\noindent\hspace*{1.2em}\kbvlabel{ಉದಾ}\hspace{0.6em}ಸರ್ಕಸ್ಸಿನಲ್ಲಿ ಪ್ರಾಣಿಗಳೂ ಜನರೂ ಸಾಹಸದ ಕೆಲಸಗಳನ್ನು ಮಾಡಿದರು. ರಾಮನೂ ಲಕ್ಷ್ಮಣನೂ ಕಪಿ ಸೈನ್ಯವೂ ರಾವಣನ ಮೇಲೆ ಯುದ್ಧ ಮಾಡಿತು.\par

\vspace{2.5pt plus 3pt minus 0.5pt}\noindent\hspace*{1.2em}\kbvnum{{[}4{]}}\hspace{0.55em}ಒಂದೇ ಕ್ರಿಯೆಗೆ ಉತ್ತಮ ಪುರುಷ ಸರ್ವನಾಮವೂ, ಇನ್ನಿತರ ಸರ್ವನಾಮಗಳೂ ಕರ್ತೃಗಳಾಗಿದ್ದಲ್ಲಿ ಕ್ರಿಯಾಪದವು ಉತ್ತಮ ಪುರುಷದಲ್ಲೇ ಇರುವುದು.\par

\noindent\hspace*{1.2em}\kbvlabel{ಉದಾ}\hspace{0.6em}ನಾನೂ ನೀನೂ ಅವನೂ ಊರಿಗೆ ಹೋದೆವು.\par

\noindent\hspace*{2.2em}ನಾನೂ ಅವನೂ ಸೇತಿಯಿಂದ ಬಂದೆವು.\par

\noindent\hspace*{2.2em}ನಾನೂ ನೀನೂ ಕೂಡಿ ಬಾಳೋಣ.\par

\vspace{2.5pt plus 3pt minus 0.5pt}\noindent\hspace*{1.2em}\kbvnum{{[}5{]}}\hspace{0.55em}ಒಂದೇ ಕ್ರಿಯೆಗೆ ಮಧ್ಯಮ ಪುರುಷ ಸರ್ವನಾಮವೂ, ಪ್ರಥಮ ಪುರುಷ ಸರ್ವನಾಮವೂ ಕರ್ತೃಗಳಾಗಿದ್ದಲ್ಲಿ ಕ್ರಿಯಾಪದವು ಮಧ್ಯಮ ಪುರುಷದಲ್ಲೇ ಇರುವುದು.\par

\noindent\hspace*{1.2em}\kbvlabel{ಉದಾ}\hspace{0.6em}ನೀನೂ ಅವನೂ ಹಣ್ಣುಗಳನ್ನು ತಿಂದಿರಿ.\par

\noindent\hspace*{2.2em}ಅವಳೂ ನೀನೂ ಸೇರಿ ಹೊಯ್ದ ಕಿತ್ತಿರಿ.\par

\kbvexercise{ಅಭ್ಯಾಸ ಪತ್ರಿಕೆ ೧೨}

\vspace{2.5pt plus 3pt minus 0.5pt}\noindent\hspace*{1.2em}\kbvnum{{[}1{]}}\hspace{0.55em}ಉತ್ತಮ ಪುರುಷ ಪುಲ್ಲಿಂಗ ಏಕವಚನದಲ್ಲಿ “ಬರು” ಧಾತು; ವರ್ತಮಾನ ಕಾಲದ ಕ್ರಿಯಾಪದ ರೂಪವನ್ನು ತಿಳಿಸಿರಿ.\par
